This chapter evaluates the prototype described in Chapter~\ref{chap:implementation} against the functional and non-functional requirements defined in Section~\ref{sec:requirements-principles} and the research questions formulated in Chapter~\ref{chap:methodology}. The evaluation combines three complementary methods: (1) unit test results that verify individual requirements through the test traceability matrix established in Section~\ref{sec:test-traceability}, (2) quantitative experiments that measure alert volume, \gls{cpu} utilisation, and memory consumption under controlled fault scenarios, and (3) a qualitative comparison with two established monitoring tools Prometheus and Nagios that contextualises the framework's capabilities against existing practice.

The evaluation does not claim that the prototype is production-ready or that it supersedes general-purpose monitoring platforms. Prometheus and Nagios serve fundamentally different purposes and operate at different architectural scales. They are included as reference points to illustrate the behavioural differences that arise from the framework's episode-aware, dependency-coordinated design specifically in alert volume, resource consumption, and the presence or absence of root-cause analysis, suppression, and episode tracking. Where the data supports a conclusion, it is stated. Where it does not, the limitation is noted.

\subsection{Evaluation Setup}\label{sec:eval-setup}

\subsubsection{Hardware Environment}

All experiments were conducted on the following hardware:

\begin{table}[htbp]
\centering
\caption{Hardware environment}
\label{tab:hardware-env}
\begin{tabular}{p{4cm} p{9cm}}
\toprule
Component & Specification \\
\midrule
CPU & Intel Core i7-8700K @ 3.70 GHz \\
Cores & 6 cores / 12 threads \\
Cache & 12 MB L3 \\
Memory & 32 GB RAM, 8 GB swap \\
Operating System & Ubuntu 24.04.4 LTS \\
Kernel & Linux 6.17 \\
Architecture & x86\_64 \\
\bottomrule
\end{tabular}
\end{table}

This hardware exceeds the resource profile of the target industrial router platform. Consequently, absolute CPU and memory measurements are not directly transferable to the embedded deployment target. The evaluation uses this environment to establish relative comparisons between the three systems under identical conditions and to verify functional correctness. Resource behaviour on the actual router platform remains an open validation point (Section~\ref{sec:eval-limitations}).

\subsubsection{Software Environment}

The prototype (\texttt{prism-police}) was deployed in a Docker container that replicates the PRISM platform's \gls{soa} environment. The container includes:

\begin{itemize}
\item The \texttt{netconfd-pro} \gls{yang}/\gls{netconf} server for north-south configuration delivery.
\item A proxy application for east-west message bus communication via ZeroMQ.
\item The Yuma \gls{netconf} client for configuration injection and operational state queries.
\end{itemize}

Prometheus (0.32.0) and Nagios (4.5.12) were deployed in the same Docker-based environment with equivalent monitoring configurations: the same telemetry sources, the same threshold rules, and the same fault injection sequences. All three systems received identical telemetry streams during each experiment. The specific versions of Prometheus (0.32.0) and Nagios (4.5.12) were recorded during the experiments, which gives exact reproducibility of the comparison.

\textbf{Limitation:} The Docker-based deployment introduces containerisation overhead (network namespacing, filesystem layering) that affects all three systems equally. The measurements reflect containerised performance, not bare-metal performance.
\subsubsection{Measurement Methodology}

Resource consumption and alert volume were measured using operating system interfaces on the Linux-based experiment host. The measurement tools and methods differ by system due to their different deployment models:

\begin{itemize}
\item \textbf{Memory (prism-police):} Resident Set Size (RSS) was read from \texttt{/proc/\textit{PID}/status} (field \texttt{VmRSS}) once after a 10-second stabilisation period following experiment start. RSS represents the portion of the process's memory held in physical RAM. A single snapshot was taken per experiment run.
\item \textbf{Memory (Prometheus + Alertmanager):} RSS of both container main processes (Prometheus and Alertmanager) was sampled once per second for the experiment duration using \texttt{ps -o rss=}. The per-second values were summed across both processes and the reported value is the arithmetic mean across all samples.
\item \textbf{Memory (Nagios):} The \texttt{MemoryCurrent} property was read from the \texttt{nagios4} systemd service unit via \texttt{systemctl show nagios4 -p MemoryCurrent} at experiment end. This value represents the current memory consumption of the service's cgroup as reported by the kernel.
\item \textbf{CPU (prism-police):} Sampled once per second via \texttt{ps -p \textit{PID} -o \%cpu=} for the experiment duration (120--300 samples depending on scenario). The reported value is the arithmetic mean across all samples.
\item \textbf{CPU (Prometheus + Alertmanager):} Measured using \texttt{pidstat -p \textit{PIDs} 1 \textit{duration}} at one-second intervals for both container main processes. The reported value is the sum of the average CPU percentages of both processes.
\item \textbf{CPU (Nagios):} The \texttt{CPUUsageNSec} property was read from the \texttt{nagios4} systemd service unit at experiment start and end. The reported CPU usage is the delta between the two readings, representing total CPU time consumed during the measurement window.
\item \textbf{Alert volume (prism-police):} Fault episode count reported by the framework at experiment end.
\item \textbf{Alert volume (Prometheus):} Notification count from a webhook receiver (a Python HTTP server on port 9099 that increments a counter for each POST received from Alertmanager).
\item \textbf{Alert volume (Nagios):} Count of \texttt{SERVICE NOTIFICATION:} entries in \texttt{/var/log/nagios4/nagios.log} with timestamps falling within the measurement window, extracted via \texttt{awk}.
\end{itemize}

Each experiment was preceded by a 10--20 second warm-up period (10 seconds for E0--E2, 20 seconds for E3 scenarios) excluded from measurement. The experiment suites were executed via automated shell scripts (\texttt{measure\_e0.sh} through \texttt{measure\_e3\_1000.sh} for prism-police, \texttt{run\_prometheus\_suite.sh} for Prometheus, and \texttt{run\_nagios\_suite.sh} for Nagios) to ensure reproducible execution sequences.

\textbf{Out-of-memory behaviour:} The prototype does not implement explicit handling for memory allocation failures. If the operating system denies a memory allocation request, the C++ runtime throws \texttt{std::bad\_alloc}, which is unhandled and terminates the process. The bounded event queue (drop-oldest policy, capacity 1000) and the absence of persistent telemetry storage are designed to prevent unbounded memory growth under normal operation. On the target platform, the service lifecycle manager would detect the terminated process and restart it.

\subsubsection{Experiment Design}

Four experiment scenarios were defined, each targeting a different operational condition:

\begin{itemize}
\item E0 Idle baseline (120 seconds). No faults injected. All telemetry sources report healthy values. This scenario establishes the baseline resource consumption of each system when no fault processing occurs.

\item E1 Single sustained fault (300 seconds). A single service reports a sustained threshold violation for the duration of the experiment. This scenario tests the system's behaviour under a persistent, isolated fault condition.

\item E2 Cascading fault with four services and one root cause (300 seconds). One root-cause service fails, and three dependent services exhibit downstream symptoms. This scenario tests dependency-aware detection, root-cause analysis, and suppression behaviour.

\item E3 Message storm under sustained fault (180 seconds, N = 100 / 500 / 1000). A sustained fault condition is maintained while the telemetry ingestion rate is increased to N samples per second. This scenario tests the system's behaviour under high telemetry load, measuring alert volume growth, CPU utilisation scaling, and memory consumption under stress.
\end{itemize}

Each scenario was executed once per system. The measurements reported are single-run values, not averages over multiple runs. This limits the statistical confidence of the results; variability across runs is not quantified. This is acknowledged as a limitation in Section~\ref{sec:eval-limitations}.

\subsection{Quantitative Results}\label{sec:quantitative-results}

\subsubsection{Alert Volume}

Table reports the total number of alerts generated by each system during each experiment scenario.

\begin{table}[htbp]
\centering
\begin{tabular}{|c|c|c|c|c|}
\hline
Scenario & Load (samples/s) & Prometheus & Nagios & prism-police \\
\hline
E1 & - & 176 & 298 & 1 \\
E2 & - & 666 & 1,201 & 1 \\
E3 & 100 & 10,016 & 18,006 & 100 \\
E3 & 500 & 37,977 & 29,711 & 500 \\
E3 & 1000 & 37,220 & 179,073 & 1,000 \\
\hline
\end{tabular}
\caption{Alert volume across experiment scenarios}
\label{tab:alert-volume}
\end{table}

\textbf{Observations:}

In scenarios E1 and E2, prism-police generates exactly one alert. This is a direct consequence of the episode model (Section~\ref{sec:episode-model}) and the at-most-once resolution guarantee (Section~\ref{sec:at-most-once}): the first threshold violation starts a fault episode and triggers a single detection event. Subsequent samples that continue to violate the same rule within the same episode do not generate additional alerts, because the \texttt{RuleStateTracker} reports no new rising-edge transition (Section~\ref{sec:rule-state-tracking}).

Prometheus generates 176 alerts in E1 and 666 in E2. Nagios generates 298 and 1,201 respectively. Both systems re-evaluate and re-fire alerts on each evaluation cycle for as long as the violation persists, because neither maintains episode state that distinguishes a new fault occurrence from an ongoing violation. This behaviour is consistent with the gap identified in Section~\ref{sec:lit-fault-models}: existing monitoring tools treat each evaluation cycle independently.

In E2, the difference between Prometheus (666) and prism-police (1) is attributable to two factors: (i) Prometheus generates alerts for all four faulted services independently, while prism-police identifies the root cause and suppresses the three downstream symptoms; and (ii) Prometheus re-fires alerts on each evaluation cycle, while prism-police fires once per episode.

In E3, prism-police generates exactly N alerts (100, 500, 1,000). The E3 scenario injects telemetry samples from N distinct telemetry sources per second, each representing an independent service metric. Each source's first violating sample starts a new fault episode and generates one alert. Within each source's episode, subsequent samples do not generate additional alerts (consistent with the episode model). The total alert count equals N because N independent episodes are created — one per source. The at-most-once resolution guarantee ensures that resolution executes only once per episode per root cause, regardless of the total sample volume. Prometheus and Nagios generate alert counts that are multiples of the sample rate, reflecting repeated re-evaluation across all sources on every evaluation cycle without episode gating.

At N=500, Nagios (29,711) generates fewer alerts than Prometheus (37,977), while at N=1000, Nagios (179,073) generates substantially more than Prometheus (37,220). This inversion suggests that Nagios's alert generation behaviour is non-linear with respect to ingestion rate, possibly due to internal queuing or check scheduling effects. The cause of this inversion was not investigated further.

\textbf{Interpretation:}

The alert volume data demonstrates that the episode model and suppression mechanism reduce the number of alerts by two to five orders of magnitude compared to stateless monitoring tools under sustained fault conditions. This reduction is a structural consequence of the design (episode gating, rising-edge detection, root-cause-only alerting) rather than a tuning artefact. The reduction directly addresses the ``redundant alerts'' and ``cascading recovery actions'' identified as gaps in Section~\ref{sec:lit-summary}.

\subsubsection{CPU Utilisation}

Table reports the CPU utilisation (in percentage of a single core) for each system during each scenario.

\begin{table}[htbp]
\centering
\begin{tabular}{|c|c|c|c|c|}
\hline
Scenario & Input Rate (samples/s) & Prometheus (\%) & Nagios (\%) & prism-police (\%) \\
\hline
E0 & - & 0.20 & 0.98 & 0.38 \\
E1 & - & 1.80 & 4.27 & 0.50 \\
E2 & - & 4.68 & 17.32 & 0.57 \\
E3 & 100 & 58.18 & 127.68 & 11.34 \\
E3 & 500 & 228.76 & 787.40 & 56.54 \\
E3 & 1000 & 313.83 & 1,035.97 & 77.49 \\
\hline
\end{tabular}
\caption{CPU utilisation across experiment scenarios}
\label{tab:cpu-utilisation}
\end{table}

Note: Values exceeding 100\% indicate utilisation across multiple cores (e.g., 228.76\% $\approx$ 2.3 cores).

\textbf{Observations:}

At idle (E0), all three systems consume minimal CPU. Prism-police (0.38\%) consumes more than Prometheus (0.20\%) at idle, which is attributable to the 100ms polling loop that executes \texttt{processAllEvents()} even when the event queue is empty. Nagios (0.98\%) consumes the most at idle, likely due to its periodic check scheduling.

Under single-fault (E1) and cascading-fault (E2) conditions, prism-police's CPU utilisation increases marginally (0.50\% and 0.57\%), while Prometheus and Nagios show larger increases (1.80\%/4.68\% and 4.27\%/17.32\% respectively). The small increase for prism-police reflects the sample-driven processing model (Principle 1): CPU usage is proportional to the number of incoming samples, not to the number of active alerts or the complexity of the fault condition.

Under message storm conditions (E3), CPU utilisation scales with ingestion rate for all three systems. At N=1000, prism-police consumes 77.49\% of a single core, Prometheus consumes 313.83\% (approximately 3.1 cores), and Nagios consumes 1,035.97\% (approximately 10.4 cores). Prism-police's CPU consumption at N=1000 is 24.7\% of Prometheus's and 7.5\% of Nagios's.

The lower CPU utilisation of prism-police under storm conditions is attributable to two factors: (i) the bounded event queue (1000 events, drop-oldest) limits the number of events processed per cycle, preventing unbounded processing backlogs; and (ii) the episode model and suppression mechanism reduce the amount of downstream processing (root-cause analysis, resolution execution, alert generation) per sample compared to systems that re-evaluate all rules and re-fire all alerts on every cycle.

\textbf{Interpretation:}

The CPU data supports the claim that the framework's processing overhead remains bounded relative to the ingestion rate (NFR4, Principle 5). However, the absolute values are measured on a desktop-class CPU and cannot be directly extrapolated to the embedded router platform. The relative comparison prism-police consuming 4× to 13× less CPU than the reference systems under storm conditions indicates that the architectural decisions (bounded queue, episode gating, suppression) contribute to reduced processing overhead. The idle CPU consumption (0.38\%) caused by the 100\,ms polling loop is an identified optimisation point: replacing the fixed-interval polling with an event-driven wake-up mechanism (e.g., a condition variable signalled by the InputBlock upon enqueue) would reduce idle CPU consumption to near zero without affecting processing behaviour under load.

\subsubsection{Memory Utilisation}

Table reports the memory consumption (in megabytes) for each system during each scenario.

\begin{table}[htbp]
\centering
\begin{tabular}{|c|c|c|c|c|}
\hline
Scenario & Input Rate (samples/s) & Prometheus (MB) & Nagios (MB) & prism-police (MB) \\
\hline
E0 & - & 196.00 & 6.70 & 13.92 \\
E1 & - & 197.00 & 6.90 & 13.98 \\
E2 & - & 207.00 & 7.30 & 14.17 \\
E3 & 100 & 219.00 & 10.60 & 15.51 \\
E3 & 500 & 247.00 & 26.10 & 18.66 \\
E3 & 1000 & 374.00 & 41.40 & 29.37 \\
\hline
\end{tabular}
\caption{Memory utilisation across experiment scenarios}
\label{tab:memory-utilisation}
\end{table}

\textbf{Observations:}

Prometheus has the highest baseline memory consumption (196 MB at idle), which reflects its time-series database, query engine, and HTTP server infrastructure. Nagios has the lowest baseline (6.70 MB), reflecting its lightweight C-based architecture. Prism-police's baseline (13.92 MB) falls between the two.

Under increasing load (E3), all three systems show memory growth. Prometheus grows from 196 MB to 374 MB (1.9× increase from E0 to E3-1000). Nagios grows from 6.70 MB to 41.40 MB (6.2× increase). Prism-police grows from 13.92 MB to 29.37 MB (2.1× increase).

Prism-police's memory growth from E0 to E3-1000 is 15.45 MB. The bounded event queue itself does not contribute to this growth: it has a fixed capacity of 1000 elements and operates as a drop-oldest ring buffer whose memory footprint is constant regardless of ingestion rate. The observed growth is instead attributable to data structures outside the queue that scale with the number of distinct telemetry sources: sample storage (one stored sample per telemetry ID), fault records (one per active violation), suppression counters (one per suppressed telemetry), the \texttt{resolvedEpisodes} set, and the \texttt{persistentFaults} vector. In E3-1000, 1000 distinct telemetry sources are active simultaneously, each requiring its own per-source state. The \texttt{resolvedEpisodes} set and \texttt{persistentFaults} vector grow monotonically because the pruning mechanism described in Section~\ref{sec:at-most-once} is not implemented in the prototype. Bounding these structures to the number of concurrently active episodes (rather than the cumulative total) is an identified optimisation point for future work.

\textbf{Interpretation:}

The memory data shows that prism-police's memory consumption ranges from 13.92~MB (E0, idle) to 29.37~MB (E3-1000, maximum tested load). For context, a representative embedded industrial router platform provides 256~MB of RAM, of which the operating system and other services consume a portion. Prism-police's peak measured consumption of 29.37~MB represents approximately 11.5\% of a 256~MB memory budget, leaving substantial headroom for other platform services. However, the test T16 (NFR4 --- bounded resource usage) is a placeholder in the current test suite and does not verify that memory consumption remains bounded over extended operational periods (days or weeks). The 2.1\texttimes{} growth factor under storm conditions indicates that memory usage is correlated with the number of concurrently active telemetry sources. Whether this growth remains bounded over time --- particularly whether the monotonically growing \texttt{resolvedEpisodes} set and \texttt{persistentFaults} vector cause unbounded accumulation during long-duration operation --- requires soak testing not performed in this evaluation. Implementing the episode-end pruning mechanism defined in Section~\ref{sec:at-most-once} would bound these structures and is identified as a necessary optimisation for production deployment.

Figure~\ref{fig:quantitative-comparison} summarises the CPU and memory utilisation across all scenarios.

\begin{figure}[htbp]
\centering
\includegraphics[width=0.8\textwidth]{quantitative_comparison_bar_chart.png}
\caption{Bar graph quantitative comparison of CPU and memory utilisation}
\label{fig:quantitative-comparison}
\end{figure}

\subsection{Qualitative Comparison}\label{sec:qualitative-comparison}

The quantitative results in Section~\ref{sec:quantitative-results} measure observable outputs (alerts, CPU, memory) but do not capture the architectural capabilities that differentiate the three systems. This section compares the systems across seven qualitative dimensions derived from the research questions and the gaps identified in the literature review.

\begin{table}[htbp]
\centering
\footnotesize
\caption{Qualitative comparison}
\label{tab:qualitative-comparison}
\begin{tabular}{p{2.8cm} p{3.5cm} p{3.5cm} p{4.8cm}}
\toprule
\textbf{Capability} & \textbf{Prometheus} & \textbf{Nagios} & \textbf{prism-police} \\
\midrule
Root-cause analysis & Not supported. No dependency model. & Not supported. Limited host dependencies; no cross-layer causal traversal. & Supported. Dependency graph traversal identifies deepest active upstream fault (T04, T05, T23, T24, T29). \\

Suppression of dependent symptoms & Not natively supported. Requires manual inhibition rules; no transitive handling. & Limited. Host-level suppression only; no service-level propagation. & Supported. Automatic transitive suppression via reverse dependency graph (\gls{bfs}) (T06, T25, T45). \\

Persistent fault tracking & Not supported. No distinction between transient and post-resolution faults. & Not supported. Check states do not retain resolution history. & Supported. Post-resolution violations tracked and exposed via \gls{netconf} (T09, T10, T31). \\

Episode-based handling & Not supported. Evaluations are stateless; deduplication label-based. & Not supported. Flap detection exists but no episode semantics. & Supported. First-class fault episodes with at-most-once resolution (T02, T15, T34, T35). \\

Determinism & Partial. Evaluation deterministic; scrape timing and routing non-deterministic. & Partial. Checks deterministic; scheduling subject to jitter. & Deterministic by design. Timestamp-ordered, dependency-aware processing (T14, T37--T41). \\

Reconfiguration without code changes & Supported. YAML-based configuration with reload. & Supported. File-based configuration; requires service reload. & Supported. \gls{yang}/\gls{netconf}-based runtime configuration (T11, T20). \\

Dashboard / operator interface & Supported. Native UI; integrates with Grafana. & Supported. Built-in web interface with history view. & Not provided. Exposes operational state via \gls{netconf} and structured logs only. \\
\bottomrule
\end{tabular}
\end{table}

The qualitative comparison shows that prism-police provides four capabilities absent in both Prometheus and Nagios: runtime root-cause analysis with dependency traversal, automatic transitive suppression of downstream symptoms, persistent fault tracking with root-cause attribution, and episode-based recovery gating with at-most-once execution semantics. These four capabilities correspond directly to the four gaps identified in Section~\ref{sec:lit-summary}.

Prometheus and Nagios provide capabilities that prism-police does not: graphical dashboards, historical data storage, notification routing, and integration with broad ecosystems of exporters and plugins. These capabilities are outside the scope of the \gls{fdir} framework (Section~\ref{sec:scope-boundaries}) but are relevant for operational deployment. A production deployment would likely combine prism-police for \gls{fdir} coordination with an external monitoring tool for visualisation and long-term telemetry storage.

\subsection{Requirements Validation}\label{sec:requirements-validation}

This section evaluates each functional and non-functional requirement against the combined evidence from unit tests (Chapter~\ref{chap:implementation}), quantitative experiments (Section~\ref{sec:quantitative-results}), and qualitative analysis (Section~\ref{sec:qualitative-comparison}). Table~\ref{tab:req-validation} summarises the validation status and supporting evidence for each requirement. All eight functional requirements are validated. Nine of eleven non-functional requirements are validated; NFR4 (bounded resource usage) and NFR5 (timely fault handling) are partially validated due to the absence of long-duration soak testing and production-platform latency measurements, respectively.
\FloatBarrier
\begin{table}[htbp]
\centering
\caption{Requirements validation summary}
\label{tab:req-validation}
\begin{tabular}{p{2cm} p{3cm} p{8cm}}
\toprule
\textbf{Requirement} & \textbf{Status} & \textbf{Evidence} \\
\midrule
FR1 & Validated & T01, T02, T30, T32, T35--T41; E1, E2, E3 \\
FR2 & Validated & T03, T23, T24 \\
FR3 & Validated & T04, T05, T23, T24, T29; E2 \\
FR4 & Validated & T06, T15, T25, T27, T34; E2 alert count \\
FR5 & Validated & T08, T09, T31 \\
FR6 & Validated & T10 \\
FR7 & Validated & T11, T20 \\
FR8 & Validated & T12, T13, T33 \\
NFR1 & Validated & T06, T23, T24, T25; E2 \\
NFR2 & Validated & T14, T37--T41 \\
NFR3 & Validated & T15, T34; E1 \\
NFR4 & Partially validated & Memory evaluation (Table~7.3); T16 placeholder \\
NFR5 & Partially validated & T17 proxy; 100 ms polling interval \\
NFR6 & Validated & T10; logging support (T18 requires infrastructure) \\
NFR7 & Validated & T19; E3 \\
NFR8 & Validated & T11, T20 \\
NFR9 & Validated & T21 \\
NFR10 & Validated & T10; \gls{netconf} integration verified \\
NFR11 & Validated & T22, T30 \\
\bottomrule
\end{tabular}
\end{table}

\subsection{Discussion}\label{sec:eval-discussion}

\subsubsection{Addressing the Research Gap}

The literature review (Chapter~\ref{chap:literature-review-results}) identified four recurring gaps: (1) the absence of runtime fault episode tracking, (2) the lack of dependency-aware recovery coordination, (3) the missing enforcement of at-most-once recovery execution per fault episode, and (4) the reliance on computational methods unsuitable for embedded platforms. The evaluation provides evidence that the prototype addresses each gap:

Gap (1) is addressed by the episode model. The E1 scenario demonstrates that a 300-second sustained fault generates exactly one alert, compared to 176 (Prometheus) and 298 (Nagios). This reduction is a direct consequence of the rising-edge detection (Section~\ref{sec:rule-state-tracking}) and episode lifecycle management (Section~\ref{sec:episode-lifecycle}) that group repeated observations into a single episode. Unit tests T02, T15, T34, and T35 verify the episode lifecycle under various conditions.

Gap (2) is addressed by the dependency graph and root-cause analysis. The E2 scenario demonstrates that a four-service cascading fault generates one alert (for the root cause) rather than four independent alerts. Unit tests T04, T23, T24, and T29 verify that root-cause traversal correctly identifies the deepest active upstream fault across multi-layer dependency chains.

Gap (3) is addressed by the resolved-episodes set. Unit tests T15 and T34 verify that resolution executes at most once per episode and that new episodes permit re-execution. The E1 alert count (1) confirms that the at-most-once guarantee holds under sustained fault conditions.

Gap (4) is addressed by the rule-based, deterministic processing model. The CPU utilisation data (Table~7.2) shows that prism-police consumes 77.49\% of a single core at 1000 samples per second, compared to 313.83\% (Prometheus) and 1,035.97\% (Nagios). The framework uses no \gls{ml}-based inference, no statistical models, and no adaptive algorithms. All processing is deterministic and reproducible (T14). The determinism claim applies to decision logic and processing order: identical inputs produce identical fault detection, root-cause identification, and resolution outcomes. The evaluation measures average CPU utilisation but does not measure worst-case execution time (\gls{wcet}) or latency jitter. Execution time analysis on the target embedded platform remains future work (Section8.4).

\subsubsection{Alert Reduction as a Structural Property}

The alert reduction observed in the experiments is not a parameter-tuning result. It is a structural consequence of three design decisions: (i) rising-edge detection ensures that only new violations generate alerts, not repeated observations; (ii) episode gating ensures that resolution executes at most once per episode; and (iii) transitive suppression ensures that downstream symptoms do not generate independent alerts while the root cause is being resolved.

This structural property means that the alert reduction holds regardless of the specific threshold values, the number of services, or the telemetry ingestion rate. The E3 experiments confirm this: at N=1000, prism-police generates 1,000 alerts (one per sample event processed) while Prometheus generates 37,220 and Nagios generates 179,073. The ratio between prism-police and the reference systems increases with ingestion rate, because the reference systems' alert counts grow with both the number of evaluation cycles and the number of active violations, while prism-police's alert count grows only with the number of new fault episodes.

\subsubsection{Resource Consumption in Context}

The CPU and memory measurements must be interpreted with two caveats. First, the evaluation hardware (i7-8700K, 32 GB \gls{ram}) is not representative of the target embedded platform. The absolute values are meaningful only for relative comparison between the three systems. Second, the measurements are single-run values without statistical repetition.

Within these constraints, the data shows that prism-police's resource consumption scales more favourably than the reference systems under increasing load. The CPU utilisation ratio (prism-police / Prometheus) decreases from 1.9 at idle (0.38/0.20) to 0.25 at E3-1000 (77.49/313.83), indicating that the framework's processing overhead grows more slowly than Prometheus's under storm conditions. This is consistent with the bounded event queue and episode-gating mechanisms that limit per-cycle processing regardless of ingestion rate.

Memory consumption for prism-police (13.92--29.37 MB) is substantially lower than Prometheus (196--374 MB) but higher than Nagios (6.70--41.40 MB) at low load. At E3-1000, prism-police (29.37 MB) and Nagios (41.40 MB) converge, while Prometheus (374 MB) remains an order of magnitude higher. The difference reflects Prometheus's time-series database, which stores historical samples, versus prism-police's sample-driven model, which does not maintain long-term telemetry archives.

\subsubsection{Limitations of the Comparison}

Prometheus and Nagios are general-purpose monitoring platforms designed for broad observability across heterogeneous infrastructure. They provide capabilities dashboards, historical storage, notification routing, plugin ecosystems that are outside the scope of the \gls{fdir} framework. The comparison is not intended to argue that prism-police replaces these tools. Rather, it demonstrates that the specific capabilities required for stateful \gls{fdir} in industrial routers (episode tracking, dependency-aware root-cause analysis, transitive suppression, at-most-once recovery) are absent in general-purpose monitoring tools and that their absence results in measurably higher alert volumes and resource consumption under fault conditions.

A fair operational comparison would require deploying prism-police alongside a monitoring tool (e.g., Prometheus for visualisation and historical storage) and measuring the combined system's behaviour. This combined deployment was not evaluated.

\subsubsection{Open Validation Gaps: NFR4 and NFR5}

Two non-functional requirements remain partially validated, and this represents the primary open validation gap of the thesis.

As discussed in Section~\ref{sec:at-most-once}, entries for ended episodes can be safely removed because episode IDs are monotonically increasing and never reused; implementing this pruning would bound the set's size to the number of concurrently active episodes rather than the cumulative total.

NFR5 (timely fault handling) is supported by the 100ms polling interval, which establishes a theoretical detection latency floor, and by the proxy measurement in T17 (below 1 second in the unit test environment). However, end-to-end latency from telemetry sample publication on the message bus through \gls{fdir} processing to resolution script completion has not been measured on the target platform under production load. The 100ms polling interval, message bus delivery latency, \gls{netconf} processing overhead, and fork+exec resolution cost may combine to produce detection-to-resolution latencies that exceed acceptable bounds on resource-constrained hardware. This measurement requires deployment on the actual router platform and is identified as future work.

\subsection{Research Question Evaluation}\label{sec:rq-evaluation}

This section maps the evaluation evidence to the five research questions defined in Chapter~\ref{chap:methodology}.

RQ1: How can faults in \gls{soa}-based industrial routers be systematically characterised to support precise detection, isolation, and recovery under embedded system constraints?

The framework characterises faults through threshold-based rule evaluation on structured telemetry samples, with rising-edge detection that distinguishes new violations from ongoing observations. The evaluation demonstrates correct detection across single-fault (E1), cascading-fault (E2), and storm (E3) scenarios. Unit tests T01, T02, T30, T32, and T35--T41 verify detection correctness under various conditions. The characterisation is limited to single-threshold comparisons on float values; multi-condition rules and graduated severity are not implemented (Section~\ref{sec:impl-limitations}). The conceptual design defines a six-level severity model (Section~\ref{sec:fault-record-creation}), but the prototype implements a simplified uniform-severity assignment where all detected faults receive \texttt{Severity::Critical}. Graduated severity assignment based on threshold proximity, violation duration, or fault type as proposed in the recovery-oriented taxonomy of Vitucci et al.~\cite{ref1,ref2} is designed as an extension point but not realised in the current prototype. The evaluation therefore validates fault detection and episode tracking but does not validate severity-differentiated recovery selection.

RQ2: How can telemetry in industrial routers be organised into a coherent model that enables stateful fault tracking across time and supports root cause localisation under bounded resource constraints?

The telemetry model organises data into definitions, rules, and samples, delivered through \gls{yang}/\gls{netconf} and the message bus respectively. The bounded event queue (1000 events, drop-oldest) and sample-driven processing ensure that telemetry processing does not exceed resource bounds. Memory consumption remains between 13.92 and 29.37 MB across all tested scenarios (Table~7.3). The model supports stateful fault tracking through the episode lifecycle and persistent fault detection. Long-duration bounded-resource verification remains open (NFR4 partially validated).

RQ3: How should service dependencies in industrial routers be represented to enable deterministic reasoning about failure propagation and to identify causal relationships between service malfunctions?

Dependencies are represented as a directed graph derived from rule declarations, with forward edges for root-cause traversal and reverse edges for transitive impact computation. The evaluation demonstrates correct root-cause identification in three-layer chains (T04), four-layer chains (T23), partial chains (T24), cyclic graphs (T05), and diamond topologies (diamond tests). The dependency model is static (declared in configuration) and does not support runtime discovery. All dependency types are treated uniformly without distinguishing mandatory from optional dependencies (Section~\ref{sec:dependency-graph-structure}).

RQ4: Which deterministic recovery strategies are suitable for industrial routers, and how can episode-aware recovery coordination prevent redundant recovery actions, oscillations, and collateral disruption?

The framework implements external script execution as the recovery strategy, with episode-aware gating (at-most-once per episode), root-cause-only resolution, and transitive suppression of downstream impacts. The E1 and E2 scenarios demonstrate that recovery executes once regardless of fault duration or the number of affected services. Unit tests T06, T15, T21, T25, T27, and T34 verify the recovery coordination mechanisms. The framework does not implement recovery retry, alternative strategies, or graduated recovery (Section~\ref{sec:impl-limitations}).

RQ5: How can fault characterisation, telemetry modelling, dependency representation, and recovery strategies be integrated into a stateful \gls{fdir} architecture that coordinates root cause recovery and suppresses dependent symptom faults?

The integrated architecture composes seven components into a seven-phase processing pipeline that unifies detection, dependency reasoning, episode tracking, and recovery coordination. The E2 scenario provides end-to-end evidence: a single root-cause fault propagates to three dependent services; the framework detects all four faults, identifies the root cause, executes resolution once, suppresses the three symptoms, and records a persistent fault if the root cause persists. Test T45 verifies this complete pipeline with a four-layer chain under realistic multi-rate sampling conditions. The qualitative comparison (Table~7.4) confirms that this integration is absent in both Prometheus and Nagios.

\subsection{Threats to Validity}\label{sec:threats-validity}

\subsubsection{Internal Validity}

The experiments use a single execution per scenario without statistical repetition. Variability due to system scheduling, container overhead, or background processes is not quantified. The measurements should be treated as indicative rather than statistically robust.

The unit test suite relies on a test-specific reset mechanism to ensure isolation between test cases: internal identifier counters are reset before each test so that test outcomes do not depend on execution order. This reset mechanism is compiled only into the test build through conditional compilation directives; the production build does not include it. Consequently, the tested binary differs slightly from the production binary in that it exposes additional internal accessors used exclusively by the test harness. The functional behaviour rule evaluation, episode management, dependency traversal, resolution execution is identical between both builds, but the presence of additional accessor methods means that the object memory layout may differ marginally. This difference does not affect the validity of functional test results but means that memory consumption measurements taken on the test build are not directly comparable to the production build.

\subsubsection{External Validity}

The evaluation hardware (desktop-class i7, 32 GB \gls{ram}) does not represent the target embedded platform. Resource consumption measurements cannot be directly transferred to industrial router hardware. The Docker-based deployment introduces containerisation overhead that affects all three systems but may not reflect bare-metal deployment characteristics.

The fault scenarios (E0--E3) are synthetic. They test specific fault patterns (single fault, cascading fault, message storm) but do not cover all fault types that occur in production industrial router deployments (e.g., intermittent faults, Byzantine faults, partial degradation without threshold crossing).

\subsubsection{Construct Validity}

The comparison with Prometheus and Nagios measures alert volume, CPU, and memory, but does not measure operational outcomes such as mean time to repair (\gls{mttr}), operator cognitive load, or false-positive rates. These metrics would require human-in-the-loop evaluation that was not performed.

The alert count metric for prism-police (1 in E1 and E2) reflects the episode model's design: one alert per episode. Whether this level of alert reduction is operationally appropriate --- or whether operators require more granular notification --- is a deployment-specific question not addressed by this evaluation.

\subsection{Limitations}\label{sec:eval-limitations}

\begin{enumerate}
\item No production-platform evaluation. All experiments were conducted on desktop-class hardware in Docker containers. The Docker-based environment used for evaluation is the same environment used for development and simulation of the PRISM platform; it replicates the \gls{soa} service structure and message bus topology but does not reproduce the hardware constraints of the target router. Resource consumption on the actual Belden/Hirschmann PRISM router platform was not measured. The 100\,ms polling interval, fork+exec resolution overhead, and memory consumption may behave differently on embedded ARM or MIPS processors with limited cache and memory bandwidth. Additionally, the evaluation host does not carry production network traffic; on the target router, CPU and memory are shared with software-based packet processing (e.g., WLAN forwarding), which competes for the same resources and may affect FDIR processing latency and memory availability under real operating conditions.

\item NFR4 not fully validated. Bounded resource usage over extended operational periods (days, weeks) was not tested. The placeholder test T16 does not perform actual memory profiling. Whether the \texttt{resolvedEpisodes} set, \texttt{persistentFaults} vector, or other tracking structures grow unboundedly over long operation remains unverified. NFR5 not validated under production conditions. Detection-to-resolution latency was measured only as a proxy in the unit test environment. The 100ms polling interval establishes a theoretical floor, but actual latency including message bus delivery, \gls{netconf} processing, and resolution script execution was not measured end-to-end on the target platform.

\item The \texttt{resolvedEpisodes} set and \texttt{persistentFaults} vector grow monotonically in the prototype because the pruning mechanism described in Section~\ref{sec:at-most-once} is not implemented.

\item No dashboard or operator interface. The prototype exposes persistent faults via \gls{netconf} but does not provide a graphical interface. The qualitative comparison (Table~7.4) identifies this as a gap relative to Prometheus and Nagios.

\item Synthetic fault scenarios only. The evaluation uses controlled, injected faults. Real-world fault patterns intermittent faults, gradual degradation, correlated multi-root-cause failures, Byzantine behaviour were not tested.

\item No combined deployment evaluation. The operational scenario of deploying prism-police alongside a monitoring tool (e.g., Prometheus for visualisation) was not tested.
\end{enumerate}

\subsection{Summary}\label{sec:eval-summary}

This chapter evaluated the stateful \gls{fdir} prototype against the requirements defined in Section~\ref{sec:requirements-principles} and the research questions formulated in Chapter~\ref{chap:methodology}, using unit test results, quantitative experiments, and a qualitative comparison with Prometheus and Nagios.

The unit test suite (47 tests plus diamond dependency tests) validates all eight functional requirements and nine of eleven non-functional requirements: resource measurements are within expected bounds for all tested scenarios, but long-duration soak testing and production-platform latency measurements were not performed.

The quantitative experiments demonstrate three measurable outcomes. First, the episode model and suppression mechanism reduce alert volume by two to five orders of magnitude compared to Prometheus and Nagios under sustained fault conditions (1 alert vs. 176--298 in E1; 1 vs. 666--1,201 in E2). Second, CPU utilisation under storm conditions (E3-1000) is 77.49\% of a single core for prism-police, compared to 313.83\% for Prometheus and 1,035.97\% for Nagios. Third, memory consumption remains between 13.92 and 29.37 MB across all scenarios, compared to 196--374 MB for Prometheus and 6.70--41.40 MB for Nagios.

The qualitative comparison identifies four capabilities present in prism-police but absent in both reference systems: runtime root-cause analysis with dependency traversal, automatic transitive suppression of downstream symptoms, persistent fault tracking with root-cause attribution, and episode-based recovery gating with at-most-once execution semantics. These four capabilities correspond directly to the four gaps identified in the literature review (Section~\ref{sec:lit-summary}).

The evaluation is limited by the use of desktop-class hardware, single-run measurements, synthetic fault scenarios, and the absence of production-platform validation. These limitations define the boundary of the current evaluation and identify directions for future validation work.